\documentclass[uplatex,a4paper,11pt,dvipdfmx]{jsarticle}

\usepackage{graphicx}
\usepackage{booktabs}
\usepackage{amsmath}
\usepackage{cite}
\usepackage{hyperref}
\usepackage{array}
\usepackage{tabularx}
\usepackage{url}

\title{LaTeX テンプレート}
\author{著者名}
\date{\today}

\begin{document}
\maketitle

\begin{abstract}
This is the abstract. 本テンプレートは up\LaTeX{} を使用しており、日本語と英語の両方を単一のファイル内で混在して執筆できます。
\end{abstract}

\section{はじめに}
本テンプレートは、日本語・英語混在の文書作成に対応しています。
図（\texttt{figs/}）、表、数式、参考文献引用の基本構成を含んでいます。

\section{基本機能}

\subsection{数式}
\begin{equation}
  E = mc^2
\end{equation}

\subsection{表（Table）}
\begin{table}[htbp]
  \centering
  \caption{サンプルの表}
  \label{tab:sample}
  \begin{tabular}{llr}
    \toprule
    項目 & 説明 & 値 \\
    \midrule
    Item A & First parameter  & 10.5 \\
    Item B & Second parameter & 20.0 \\
    \bottomrule
  \end{tabular}
\end{table}

\subsection{図（Figure）}
画像ファイルは \texttt{figs/} ディレクトリに配置して読み込みます。
% \begin{figure}[htbp]
%   \centering
%   \includegraphics[width=0.7\linewidth]{figs/example.png}
%   \caption{サンプルの図}
%   \label{fig:sample}
% \end{figure}

\subsection{参考文献の引用}
参考文献は \texttt{references.bib} に記述し、\texttt{\textbackslash cite} コマンドで引用します \cite{example2024}。

\bibliographystyle{jplain}
\bibliography{references}

\end{document}

